% Standalone section for inclusion in the parent report.
% Required packages: amsmath, amssymb, amsthm, hyperref.
% The parent document supplies theorem, lemma, corollary, and remark environments.
% All proofs below are written independently; no part is a claimed Lean formalization.

\section{Conditioning a spectral gamma sum to be small}
\label{mc:sec:general}

The small-deviation result in the supplied article determines the probability
that a reciprocal-completed-$L$ random variable is very small.  We determine
the allocation of that small value among its spectral coordinates.  The
result has three scales: a deterministic allocation profile, Gaussian
fluctuations constrained by their total, and an exponential distance from
the conditioning boundary.  A second theorem gives a sharp threshold for
how many low spectral coordinates can simultaneously behave as independent
gamma variables.

The probabilistic statements hold for every regularly varying spectrum.
The role of RH or Dirichlet GRH in the application is to identify the
resulting probability law with the reciprocal of a completed $L$-function.
The general statements below require neither arithmetic input nor simple
spectral points.

\subsection{Definitions and main statement}

Let
\[
 0<\gamma_1\leq\gamma_2\leq\cdots,\qquad \gamma_j\longrightarrow\infty,
 \qquad N(T)=\#\{j:\gamma_j\leq T\},
\]
where repeated values are retained.  Assume that $N$ is regularly varying
with an index $\alpha>0$, meaning
\begin{equation}\label{mc:eq:RV}
 \frac{N(rT)}{N(T)}\longrightarrow r^\alpha
 \quad\text{for every }r>0\quad(T\longrightarrow\infty).
\end{equation}
Thus one may write $N(T)=T^\alpha L(T)$, with $L$ slowly varying;
a fixed positive constant may be absorbed into $L$.

Fix $p>\alpha$ and $\theta>0$.  Let $G_j$ be independent gamma variables
with shape $\theta$ and rate one, so their density is
$z^{\theta-1}e^{-z}/\Gamma(\theta)$ on $z>0$.  Define
\begin{equation}\label{mc:eq:XK}
 X=\sum_{j\geq1}\frac{G_j}{\gamma_j^p},\qquad
 K(t)=\theta\sum_{j\geq1}\log\left(1+\frac{t}{\gamma_j^p}\right).
\end{equation}
The series for $X$ converges almost surely and in $L^1$; in particular,
$\mathbb E e^{-tX}=e^{-K(t)}$ for $t\geq0$.

For $0<x<\mathbb E X$, let $t=t_x$ be the unique positive solution of
\begin{equation}\label{mc:eq:saddle}
 K'(t)=\theta\sum_j\frac{1}{\gamma_j^p+t}=x,
 \qquad u=t^{1/p}.
\end{equation}
Throughout, $x\downarrow0$, and hence $t,u\longrightarrow\infty$.
For $c\in[0,\infty]$, put
\begin{align}
 X_x(c)&=\sum_{\gamma_j\leq cu}\frac{G_j}{\gamma_j^p},
 &m_x(c)&=\theta\sum_{\gamma_j\leq cu}\frac{1}{\gamma_j^p+t},
 \label{mc:eq:partials}\\
 \sigma_x^2&=\theta\sum_j\frac{1}{(\gamma_j^p+t)^2},
 &Y_x&=t(x-X).
 \label{mc:eq:sigma}
\end{align}
At $c=\infty$ the sums run over every $j$, so $X_x(\infty)=X$ and
$m_x(\infty)=x$.  The deterministic centering in
\eqref{mc:eq:partials} is the \emph{exact tilted mean}; replacing it by only
its leading asymptotic is generally invalid on the fluctuation scale.

Write
\begin{align}
 C_1&=\int_0^\infty\frac{\alpha r^{\alpha-1}}{1+r^p}\,dr
     =\frac{\alpha\pi}{p\sin(\pi\alpha/p)},
 \label{mc:eq:C1}\\
 C_2&=\int_0^\infty\frac{\alpha r^{\alpha-1}}{(1+r^p)^2}\,dr
     =\left(1-\frac\alpha p\right)C_1,
 \label{mc:eq:C2}\\
 A(c)&=\frac1{C_1}\int_0^c\frac{\alpha r^{\alpha-1}}{1+r^p}\,dr,
 &V(c)&=\frac1{C_2}\int_0^c\frac{\alpha r^{\alpha-1}}{(1+r^p)^2}\,dr.
 \label{mc:eq:profiles}
\end{align}
Both $A$ and $V$ increase continuously from zero to one.  Let $B$ be a
standard Brownian bridge, that is, a centered Gaussian process with
\[
 \operatorname{Cov}(B(s),B(r))=\min(s,r)-sr,\qquad 0\leq s,r\leq1.
\]
For probability measures $P,Q$ on the same measurable space, use the
normalization
$d_{\mathrm{TV}}(P,Q)=\sup_E|P(E)-Q(E)|$; for densities this is one half
of their $L^1$ distance.

\begin{theorem}[Joint conditional spectral law]\label{mc:thm:joint}
Under the preceding assumptions:
\begin{enumerate}
\item The exact saddle and fluctuation size satisfy
\begin{equation}\label{mc:eq:scales}
 tx\sim\theta C_1N(u),\qquad
 t^2\sigma_x^2\sim\theta C_2N(u),\qquad
 \frac{\sigma_x}{x}\sim
 \frac{\sqrt{C_2}}{\sqrt\theta C_1\sqrt{N(u)}}.
\end{equation}

\item Fix $0<c_1<\cdots<c_k<\infty$ and an integer $m\geq0$.
Conditional on $X\leq x$, the vector
\begin{equation}\label{mc:eq:jointvector}
 \left(
 \left(\frac{X_x(c_i)-m_x(c_i)}{\sigma_x}\right)_{i=1}^k,
 Y_x,
 \left(\frac{tG_j}{\gamma_j^p}\right)_{j=1}^m
 \right)
\end{equation}
converges in total variation to
\begin{equation}\label{mc:eq:jointlimit}
 \left((B(V(c_i)))_{i=1}^k,\ E,\ (\Gamma_j)_{j=1}^m\right),
\end{equation}
where $E$ has the exponential distribution of rate one, the $\Gamma_j$
are independent gamma variables of shape $\theta$ and rate one, and these
three objects are mutually independent.

\item The allocation profile holds uniformly in the cutoff:
\begin{equation}\label{mc:eq:uniformLLN}
 \sup_{0\leq c\leq\infty}
 \left|\frac{X_x(c)}{x}-A(c)\right|
 \longrightarrow0
 \quad\text{in probability conditional on }X\leq x.
\end{equation}
The same conclusion holds with $X$ in place of $x$ in the denominator.
\end{enumerate}
\end{theorem}

The total-variation assertion is finite dimensional.  It is deliberately
stated only at finite positive cutoffs: the standardized value at
$c=\infty$ equals $-Y_x/(t\sigma_x)$, and a continuously distributed
variable cannot converge in total variation to the constant zero.  No
functional central limit theorem in a path space is asserted here.

\subsection{Summation over a regularly varying spectrum}

\begin{lemma}[Spectral sums]\label{mc:lem:sums}
If $q>\alpha/p$, then
\begin{equation}\label{mc:eq:spectral-sum}
 \frac1{N(u)}\sum_j\frac1{(1+(\gamma_j/u)^p)^q}
 \longrightarrow
 C_q:=\int_0^\infty\frac{\alpha r^{\alpha-1}}{(1+r^p)^q}\,dr.
\end{equation}
The analogous convergence holds when the sum and integral are restricted
to a fixed interval $(a,b]$ with $0\leq a<b\leq\infty$.
Moreover, $\sum_j\gamma_j^{-p}<\infty$, and, as $h\longrightarrow\infty$,
\begin{equation}\label{mc:eq:power-sum}
 \sum_{\gamma_j\leq h}\gamma_j^{2p}
 \sim\frac{\alpha}{\alpha+2p}\,h^{2p}N(h).
\end{equation}
\end{lemma}

\begin{proof}
The measures $N(u)^{-1}\sum_j\delta_{\gamma_j/u}$ converge on every
compact subinterval of $(0,\infty)$ to the measure
$\alpha r^{\alpha-1}\,dr$, by \eqref{mc:eq:RV} and approximation of
continuous functions by step functions.  There is no difficulty at zero:
the integrand in \eqref{mc:eq:spectral-sum} is at most one, and the mass
below $\delta$ has limit $\delta^\alpha$.

For completeness, the required uniform control at infinity follows
directly from dyadic ratios.  Given $\varepsilon>0$, regular variation
implies
\[
 \frac{N(2T)}{N(T)}\leq2^{\alpha+\varepsilon}
 \quad(T\geq T_0).
\]
Iterating and using monotonicity gives, for $u\geq T_0$ and $r\geq1$,
$N(ru)/N(u)\leq C_\varepsilon r^{\alpha+\varepsilon}$.
Choose $0<\varepsilon<pq-\alpha$.  The part with $\gamma_j/u>R\geq1$
is at most
\[
 C\sum_{\ell\geq0}(2^\ell R)^{-pq}
                  (2^{\ell+1}R)^{\alpha+\varepsilon},
\]
which tends to zero as $R\to\infty$, uniformly in large $u$.
This proves the sum limit and its interval version.  The same dyadic
estimate, with an exponent between $\alpha$ and $p$, proves
$\sum_j\gamma_j^{-p}<\infty$.

Finally, Stieltjes integration by parts gives
\[
 \sum_{\gamma_j\leq h}\gamma_j^{2p}
 =h^{2p}N(h)-2p\int_0^h r^{2p-1}N(r)\,dr.
\]
After division by $h^{2p}N(h)$, dominated convergence applies to the
integral on $[0,1]$: $N(hs)/N(h)\leq1$ there and tends to $s^\alpha$.
The resulting limit is
$1-2p/(\alpha+2p)=\alpha/(\alpha+2p)$.
The constants \eqref{mc:eq:C1}--\eqref{mc:eq:C2} follow from the beta
integral after substituting $s=r^p$.
\end{proof}

\subsection{Tilting and a uniform local limit}

Since $\mathbb EX=\theta\sum_j\gamma_j^{-p}<\infty$, the nonnegative
series defining $X$ is finite almost surely.  Its Laplace transform in
\eqref{mc:eq:XK} follows by monotone convergence of the partial sums and
the gamma Laplace transform.  Termwise differentiation is justified on
every compact set of $t\geq0$.  The function $K'$ is continuous, strictly
decreasing from $\mathbb EX$ to zero, proving existence and uniqueness of
\eqref{mc:eq:saddle}.  The event $X\leq x$ has positive probability for
every $x>0$: choose a finite initial segment with expected remaining tail
less than $x/4$, apply Markov's inequality to the tail, and require each
of the finitely many initial summands to be sufficiently small.

Introduce the tilted probability measure
\begin{equation}\label{mc:eq:tilt}
 d\mathbb P_t=e^{-tX+K(t)}\,d\mathbb P.
\end{equation}
Under $\mathbb P_t$, the variables
$\widetilde G_j=(1+t/\gamma_j^p)G_j$ are independent gamma variables
of shape $\theta$ and rate one.  Indeed, the assertion follows first for
every finite group of coordinates by integrating out its independent
tail, and those finite-dimensional identities characterize the product
law.  Set
\begin{equation}\label{mc:eq:av}
 a_j=\frac{t}{\gamma_j^p+t}
     =\frac1{1+(\gamma_j/u)^p},\qquad
 v=\theta\sum_j a_j^2=t^2\sigma_x^2,
 \qquad S=tX=\sum_j a_j\widetilde G_j.
\end{equation}
The mean of $S$ under $\mathbb P_t$ is $tx$, and its variance is $v$.
Lemma~\ref{mc:lem:sums} gives \eqref{mc:eq:scales} and $v\to\infty$.

\begin{lemma}[Uniform local Gaussian limit on a spectral band]
\label{mc:lem:local}
Let $I=(a,b]\subset(0,\infty)$, with $0\leq a<b\leq\infty$.
Delete any fixed finite set of indices if desired, and write
\[
 Z_{I,u}=\frac1{\sqrt v}
 \sum_{\gamma_j/u\in I}a_j(\widetilde G_j-\theta).
\]
For all sufficiently large $u$, this variable has a continuous density
$f_{I,u}$, and
\begin{equation}\label{mc:eq:local}
 \sup_z|f_{I,u}(z)-\varphi_{\lambda_I}(z)|\longrightarrow0,
 \qquad
 \int_{\mathbb R}|f_{I,u}-\varphi_{\lambda_I}|\longrightarrow0,
\end{equation}
where $\varphi_\lambda$ is the centered normal density of variance
\[
 \lambda_I=\frac1{C_2}
 \int_I\frac{\alpha r^{\alpha-1}}{(1+r^p)^2}\,dr>0.
\]
The same conclusion holds for the whole spectrum with any lower block
$\{\gamma_j\leq h(u)\}$ removed, provided $h(u)=o(u)$; its limiting
variance is then one.
\end{lemma}

\begin{proof}
For a band, Lemma~\ref{mc:lem:sums} gives
$\theta\sum_Ia_j^2/v\to\lambda_I$.  For a fixed Fourier variable $z$,
the logarithm of the centered characteristic function is
\begin{align*}
 \log\phi_{I,u}(z)
 &=\theta\sum_I\left[-\frac{iza_j}{\sqrt v}
              -\log\left(1-\frac{iza_j}{\sqrt v}\right)\right]\\
 &=-\frac{\theta z^2}{2v}\sum_Ia_j^2
     +O_\theta\left(\frac{|z|^3\sum_Ia_j^3}{v^{3/2}}\right)
 \longrightarrow-\lambda_Iz^2/2.
\end{align*}
Here $0<a_j\leq1$, so the remainder is $O_z(v^{-1/2})$.

We now prove the Fourier integrability needed for a local, rather than
merely weak, limit.  Choose a compact interval $[a',b']$ strictly inside
$I$ with $0<a'<b'<\infty$.  There are at least $dN(u)$ spectral indices
in this interval for all large $u$, each with
$a_j\geq(1+(b')^p)^{-1}$.  Since $v\asymp N(u)$,
\begin{equation}\label{mc:eq:fourier-envelope}
 |\phi_{I,u}(z)|
 =\prod_I\left(1+\frac{z^2a_j^2}{v}\right)^{-\theta/2}
 \leq\left(1+\frac{Az^2}{N(u)}\right)^{-DN(u)}
\end{equation}
for constants $A,D>0$.  A finite deletion does not change this estimate.
Write $n=N(u)$.  For $|z|\leq\sqrt n$, the right side is at most
$\exp(-ADz^2/(1+A))$.  For the remaining tail,
\begin{align*}
 \int_{\sqrt n}^\infty
     \left(1+\frac{Az^2}{n}\right)^{-Dn}\,dz
 &=\sqrt n\int_1^\infty(1+Ar^2)^{-Dn}\,dr\\
 &\leq C\sqrt n(1+A)^{-Dn/2}\longrightarrow0,
\end{align*}
where, for large $n$, the remaining half of the power is bounded by
$(1+Ar^2)^{-1}$.  These estimates show uniform integrability of the
characteristic functions and convergence in $L^1$ to
$e^{-\lambda_Iz^2/2}$.  Fourier inversion proves uniform convergence of
the densities.  Their $L^1$ convergence follows from pointwise
convergence and equality of their integrals to one (Scheff\'e's lemma).

For a deleted lower block, $N(h)/N(u)\to0$.  To see this without any
uniform regular-variation theorem, fix $\delta>0$ and use
$h\leq\delta u$ eventually, so
$\limsup N(h)/N(u)\leq\delta^\alpha$, then let $\delta\downarrow0$.
The removed variance is at most $\theta N(h)=o(v)$.  The interval
$[u,2u]$ remains available for the Fourier envelope, and the same proof
applies with limiting variance one.
\end{proof}

\subsection{Proof of the joint conditional law}

\begin{proof}[Proof of Theorem~\ref{mc:thm:joint}]
Only the conditional conclusions remain.  Under $\mathbb P_t$, put
\[
 W=\frac{tX-tx}{\sqrt v},\qquad
 D_u=\mathbb E_t\bigl[e^{\sqrt vW}\mathbf1_{\{W\leq0\}}\bigr].
\]
The density $f_u$ of $W$ satisfies Lemma~\ref{mc:lem:local} with
limiting variance one.  On writing $y=-\sqrt vW$, we obtain
\begin{equation}\label{mc:eq:D}
 \sqrt vD_u=\int_0^\infty e^{-y}f_u(-y/\sqrt v)\,dy
 \longrightarrow\varphi_1(0)=\frac1{\sqrt{2\pi}}.
\end{equation}
The density of $\mathbb P(\,\cdot\mid X\leq x)$ with respect to
$\mathbb P_t$ is exactly
\begin{equation}\label{mc:eq:condition-weight}
 D_u^{-1}e^{\sqrt vW}\mathbf1_{\{W\leq0\}}.
\end{equation}

First omit the $m$ specified low coordinates.  Partition the remaining
indices into the $k+1$ bands
\[
 I_1=(0,c_1],\ I_2=(c_1,c_2],\ldots,
 I_k=(c_{k-1},c_k],\ I_{k+1}=(c_k,\infty).
\]
For large $u$, all $m$ omitted indices belong to $I_1$.
The centered sums on these bands, divided by $\sqrt v$, have independent
densities $f_{i,u}$ from Lemma~\ref{mc:lem:local}.  Their limiting
variances are
\[
 \lambda_i=V(c_i)-V(c_{i-1})\quad(1\leq i\leq k),
 \qquad \lambda_{k+1}=1-V(c_k),\qquad c_0=0.
\]
Let $U_i$ denote the standardized increment on the \emph{full} band
$I_i$, including the omitted coordinates if $i=1$.  For specified values
$r=(r_1,\ldots,r_m)>0$ of the scaled low coordinates
$tG_j/\gamma_j^p=a_j\widetilde G_j$, put
\[
 \delta_u(r)=\frac1{\sqrt v}\sum_{j=1}^m(r_j-\theta a_j).
\]
Under the tilted measure their joint density is
\[
 q_u(r)=\prod_{j=1}^m
 \frac{r_j^{\theta-1}e^{-r_j/a_j}}
      {\Gamma(\theta)a_j^\theta}.
\]
For fixed $r>0$, $q_u(r)$ tends to the product $q(r)$ of unit-rate gamma
densities, and $\delta_u(r)\to0$.

Transform the last band increment to $y=-\sqrt v\sum_{i=1}^{k+1}U_i$.
The Jacobian is $1/\sqrt v$.  Equations
\eqref{mc:eq:condition-weight} and \eqref{mc:eq:D} show that the joint
conditional density of $(U_1,\ldots,U_k,Y_x,r)$ is
\begin{equation}\label{mc:eq:jointdensity}
 \frac{e^{-y}\mathbf1_{\{y\geq0\}}q_u(r)}{\sqrt vD_u}
 f_{1,u}(z_1-\delta_u(r))
 \prod_{i=2}^k f_{i,u}(z_i)
 f_{k+1,u}\left(-\frac{y}{\sqrt v}-\sum_{i=1}^kz_i\right).
\end{equation}
For almost every fixed $(z,y,r)$, this converges to
\begin{equation}\label{mc:eq:limitdensity}
 \frac{e^{-y}\mathbf1_{\{y\geq0\}}q(r)}{\varphi_1(0)}
 \prod_{i=1}^k\varphi_{\lambda_i}(z_i)
 \varphi_{\lambda_{k+1}}\left(-\sum_{i=1}^kz_i\right).
\end{equation}
The product of normal densities in \eqref{mc:eq:limitdensity}, divided
by $\varphi_1(0)$, integrates to one: it is the conditional density of
independent centered normal increments of variances $\lambda_i$ given
that their sum is zero.  Consequently the entire limiting density
integrates to one.  Scheff\'e's lemma now gives total-variation
convergence of \eqref{mc:eq:jointdensity} to
\eqref{mc:eq:limitdensity}.

The cumulative increments have covariance
\[
 \sum_{i\leq\min(a,b)}\lambda_i
 -\left(\sum_{i\leq a}\lambda_i\right)
  \left(\sum_{i\leq b}\lambda_i\right)
 =V(\min(c_a,c_b))-V(c_a)V(c_b).
\]
Their distribution is therefore that of $(B(V(c_i)))_{i=1}^k$.
An invertible linear cumulative-sum transformation preserves total
variation, and the factorization in \eqref{mc:eq:limitdensity} proves
all the independence assertions.

For the uniform law of large numbers, Lemma~\ref{mc:lem:sums} gives
$m_x(c)/x\to A(c)$ for every finite $c$, while
\eqref{mc:eq:scales} gives $\sigma_x/x\to0$.  The finite-dimensional
limit therefore implies $X_x(c)/x\to A(c)$ in conditional probability
at every finite positive $c$.  At infinity,
\[
 \frac Xx=1-\frac{Y_x}{tx}\longrightarrow1,
\]
because $Y_x$ is tight and $tx\to\infty$.  Choose a finite grid of
cutoffs, including infinity, on which consecutive values of the
continuous function $A$ differ by at most $\varepsilon$.  Monotonicity
of $X_x(c)$ brackets every intermediate value between its two grid
values.  A union bound at the finitely many endpoints proves
\eqref{mc:eq:uniformLLN}.  Finally $X/x\to1$, so division by $X$ gives
the same uniform limit.
\end{proof}

\section{A sharp threshold for simultaneous gamma equilibration}
\label{mc:sec:threshold}

The fixed-coordinate conclusion of Theorem~\ref{mc:thm:joint} does not
specify how many coordinates may be treated as independent at once.
The next theorem gives an exact asymptotic answer in total-variation
distance.

Let $h=h(u)>0$ satisfy $h=o(u)$, and define the finite lower block
\[
 J_u=\{j:\gamma_j\leq h\},\qquad
 s_j=(\gamma_j/u)^p\quad(j\in J_u),\qquad
 \tau_u^2=\theta\sum_{j\in J_u}s_j^2
          =\frac\theta{u^{2p}}
            \sum_{\gamma_j\leq h}\gamma_j^{2p}.
\]
Let $Q_u$ be the conditional law of
$(tG_j/\gamma_j^p)_{j\in J_u}$ given $X\leq x$.  On the same
finite-dimensional space put
\[
 P_u=\bigotimes_{j\in J_u}\operatorname{Gamma}(\theta,1+s_j),
 \qquad
 R_u=\bigotimes_{j\in J_u}\operatorname{Gamma}(\theta,1),
\]
where the second gamma parameter is the rate.  Empty products have
their usual one-point interpretation.

\begin{theorem}[Sharp growing-block law]\label{mc:thm:threshold}
For every $h(u)=o(u)$:
\begin{enumerate}
\item The conditioning is asymptotically equivalent, on the whole
growing block, to independent exponential tilting:
\begin{equation}\label{mc:eq:canonical}
 d_{\mathrm{TV}}(Q_u,P_u)\longrightarrow0.
\end{equation}
More strongly, the joint conditional law of the block and $Y_x$ is
at total-variation distance tending to zero from
$P_u\otimes\operatorname{Exp}(1)$.

\item If $\Phi$ is the standard normal distribution function, then
\begin{equation}\label{mc:eq:sharpTV}
 d_{\mathrm{TV}}(Q_u,R_u)
 =2\Phi(\tau_u/2)-1+o(1).
\end{equation}
In particular, convergence to independent unit-rate gamma coordinates
holds if and only if $\tau_u\to0$.  If $\tau_u\to\tau\in(0,\infty)$,
the limiting distance is $2\Phi(\tau/2)-1$; if $\tau_u\to\infty$, the
distance tends to one.  The same asymptotic distance holds when the
boundary slack is included and compared with an independent
$\operatorname{Exp}(1)$ variable.

\item When $h\to\infty$, the threshold parameter has the spectral form
\begin{equation}\label{mc:eq:tauRV}
 \tau_u^2\sim
 \frac{\theta\alpha}{\alpha+2p}\,
 N(h)\left(\frac hu\right)^{2p}.
\end{equation}
\end{enumerate}
\end{theorem}

\begin{proof}
\emph{Conditioning versus tilting.}
Work under $\mathbb P_t$.  With $S=tX$ from \eqref{mc:eq:av}, write
$S=S_B+S_R$ for the lower block and its independent remainder, and let
$\mu_B,\mu_R$ be their means.  Thus $\mu_B+\mu_R=tx$.
The distribution of the block under $\mathbb P_t$ is exactly $P_u$.
Moreover,
\[
 \operatorname{Var}_t(S_B)=\theta\sum_{J_u}a_j^2
 \leq\theta N(h)=o(v).
\]
By Lemma~\ref{mc:lem:local}, the density $g_u$ of
$(S_R-\mu_R)/\sqrt v$ converges uniformly to $\varphi_1$, and is
uniformly bounded for all large $u$.

Relative to $P_u\otimes\operatorname{Exp}(1)$, the joint conditional
law of the block values $r=(r_j)_{j\in J_u}$ and the boundary slack $y$
has the exact density
\begin{equation}\label{mc:eq:blockratio}
 L_u(r,y)=\frac1{\sqrt vD_u}
 g_u\left(-\frac{\sum_{j\in J_u}r_j-\mu_B+y}{\sqrt v}\right).
\end{equation}
Indeed, after fixing the block, the remainder must have value
$tx-\sum r_j-y$, and the same change of variables as before supplies
the Jacobian $1/\sqrt v$ and the factor $e^{-y}$.
Under the reference measure in \eqref{mc:eq:blockratio},
\[
 \frac{\sum r_j-\mu_B+y}{\sqrt v}\longrightarrow0
 \quad\text{in probability},
\]
by the variance estimate and the exponential law of $y$.
Uniform convergence of $g_u$ and \eqref{mc:eq:D} imply
$L_u\to1$ in probability.  The same estimates bound $L_u$ uniformly,
so $\mathbb E|L_u-1|\to0$.  This proves the joint assertion and its
marginal \eqref{mc:eq:canonical}.

\emph{The exact statistical transition.}
Under $R_u$, let $Z_j$ be its independent unit-rate gamma coordinates.
The log likelihood of $P_u$ relative to $R_u$ is
\begin{equation}\label{mc:eq:LAN}
 \ell_u=\sum_{J_u}\{\theta\log(1+s_j)-s_jZ_j\}.
\end{equation}
Since $\max_{J_u}s_j\leq(h/u)^p\to0$,
\begin{align*}
 \mathbb E_{R_u}\ell_u
   &=-\frac{\tau_u^2}{2}
     +O\left(\max_{J_u}s_j\,\tau_u^2\right),\\
 \operatorname{Var}_{R_u}(\ell_u)&=\tau_u^2.
\end{align*}
If $\tau_u\to\tau\in(0,\infty)$, the Lyapunov estimate
\[
 \sum_{J_u}\mathbb E|s_j(Z_j-\theta)|^3
 \leq C_\theta\max_{J_u}s_j\sum_{J_u}s_j^2\longrightarrow0
\]
proves
$\ell_u\Rightarrow\mathcal N(-\tau^2/2,\tau^2)$.
For $\tau=0$, the same conclusion with a point mass at zero follows
directly from the first two moments.
The identity
\[
 d_{\mathrm{TV}}(P_u,R_u)
 =\mathbb E_{R_u}(1-e^{\ell_u})_+
\]
involves a bounded continuous function of $\ell_u$.  Hence, for finite
$\tau$, the limiting distance is
\begin{align*}
 \mathbb E(1-e^Z)_+
 &=\Phi(\tau/2)-\Phi(-\tau/2)\\
 &=2\Phi(\tau/2)-1,
 \qquad Z\sim\mathcal N(-\tau^2/2,\tau^2).
\end{align*}

If $\tau_u\to\infty$, the Hellinger affinity is explicit:
\begin{equation}\label{mc:eq:affinity}
 \int\sqrt{dP_u\,dR_u}
 =\prod_{J_u}\left(\frac{2\sqrt{1+s_j}}{2+s_j}\right)^\theta,
 \qquad
 \log\int\sqrt{dP_u\,dR_u}
 =-\frac{\tau_u^2}{8}(1+o(1)).
\end{equation}
It tends to zero.  Since
$1-d_{\mathrm{TV}}(P_u,R_u)=\int\min(dP_u,dR_u)$ is at most this
affinity, the total-variation distance tends to one.

Every sequence of $\tau_u\in[0,\infty)$ has a subsequence converging
in $[0,\infty]$.  The preceding arguments on each such subsequence
give the full asymptotic formula
$d_{\mathrm{TV}}(P_u,R_u)=2\Phi(\tau_u/2)-1+o(1)$, even when
$\tau_u$ has no limit.  The triangle inequality and
\eqref{mc:eq:canonical} give \eqref{mc:eq:sharpTV}.  Tensoring both
measures with the same exponential law preserves total variation,
which proves the statement about the slack.  Finally,
\eqref{mc:eq:tauRV} is exactly \eqref{mc:eq:power-sum}.
\end{proof}

\subsection{The larger threshold for independence with unequal rates}

Unit-rate equilibration and independence with the correct unequal rates
have different ranges.  To describe the latter, drop the restriction
$h=o(u)$ and retain the definitions of $Q_u$ and $P_u$ for an arbitrary
positive cutoff $h=h(u)$.  Define
\begin{equation}\label{mc:eq:Dprofile}
 \mathcal D(\rho)=2\left[
 \Phi\left(\sqrt{\frac{-\log(1-\rho)}{\rho}}\right)
 -\Phi\left(\sqrt{
       \frac{-(1-\rho)\log(1-\rho)}{\rho}}\right)\right],
 \qquad 0<\rho<1,
\end{equation}
and extend it continuously by $\mathcal D(0)=0$ and
$\mathcal D(1)=1$.

\begin{theorem}[Sharp boundary for independent tilted coordinates]
\label{mc:thm:macroscopic}
For every positive cutoff function $h(u)$,
\begin{equation}\label{mc:eq:macroTV}
 d_{\mathrm{TV}}(Q_u,P_u)
 =\mathcal D\bigl(V(h/u)\bigr)+o(1).
\end{equation}
Consequently, $d_{\mathrm{TV}}(Q_u,P_u)\to0$ if and only if $h=o(u)$.
If $h/u\to c\in(0,\infty)$, its limit is the strictly positive number
$\mathcal D(V(c))$; if $h/u\to\infty$, its limit is one.
\end{theorem}

\begin{proof}
Suppose first that $h/u\to c\in(0,\infty)$ and put $\rho=V(c)$.
The sum lemma and the local limit lemma apply to the moving cutoff:
the limiting variances follow by bracketing $h/u$ between $c-\varepsilon$
and $c+\varepsilon$, and fixed compact interior intervals still supply
the Fourier envelopes.  Therefore, under $P_u$, the centered block sum
\[
 Z_u=(S_B-\mu_B)/\sqrt v
\]
converges in distribution to $Z\sim\mathcal N(0,\rho)$, while the
remainder density $g_u$ converges uniformly to
$\varphi_{1-\rho}$.  Integrating \eqref{mc:eq:blockratio} over the
exponential slack shows that the marginal likelihood of $Q_u$
relative to $P_u$ differs, uniformly in the block values, by $o(1)$
from
\[
 r(Z_u)=\frac{\varphi_{1-\rho}(-Z_u)}{\varphi_1(0)}
       =(1-\rho)^{-1/2}
         \exp\left(-\frac{Z_u^2}{2(1-\rho)}\right).
\]
Uniformity follows from uniform convergence of $g_u$,
\eqref{mc:eq:D}, and the bounded derivative of the limiting normal
density: the additional shift $y/\sqrt v$ has average $1/\sqrt v$
under the exponential law.  The function $r$ is bounded and
continuous.  Thus
\[
 d_{\mathrm{TV}}(Q_u,P_u)
 \longrightarrow\tfrac12\mathbb E|r(Z)-1|.
\]
Multiplying the density of $\mathcal N(0,\rho)$ by $r$ gives the
density of $\mathcal N(0,\rho(1-\rho))$.  Their densities cross at
\[
 |z|=\sqrt{-(1-\rho)\log(1-\rho)}.
\]
Integrating the difference on the interval between these crossing
points gives precisely \eqref{mc:eq:Dprofile}.

The case $h/u\to0$ is Theorem~\ref{mc:thm:threshold}.
If $h/u\to\infty$, restrict to the smaller block with cutoff $bu$
for a fixed $b>0$.  Total variation cannot increase under this
marginalization, so
\[
 \liminf d_{\mathrm{TV}}(Q_u,P_u)\geq\mathcal D(V(b)).
\]
Letting $b\to\infty$ makes the right side tend to one.
The formula for an arbitrary $h/u$ now follows by taking subsequences
in the compactified interval $[0,\infty]$.  The claimed equivalence
uses that $V(c)>0$ and $\mathcal D(V(c))>0$ for every $c>0$.
\end{proof}

\subsection{The explicit transition for the reciprocal-$\xi$ law}

For the positive ordinates of Riemann zeros, counted with multiplicity,
$N(T)\sim T\log T/(2\pi)$.  Taking $p=2$ gives $\alpha=1$,
$C_1=\pi/2$, and $C_2=\pi/4$.  Consequently
\begin{equation}\label{mc:eq:riemannprofiles}
 A(c)=\frac2\pi\arctan c,
 \qquad
 V(c)=\frac2\pi\left(\arctan c+\frac{c}{1+c^2}\right),
\end{equation}
and, in the last equivalence below, assuming $h\to\infty$,
\begin{equation}\label{mc:eq:riemannscales}
 x\sim\frac{\theta\log u}{4u},\qquad
 \sigma_x^2\sim\frac{\theta\log u}{8u^3},\qquad
 \tau_u^2\sim\frac{\theta h^5\log h}{10\pi u^4}.
\end{equation}
These same leading formulas hold for a fixed primitive real Dirichlet
character, with central zeros omitted.

In particular, if
\begin{equation}\label{mc:eq:criticalheight}
 h=c\left(\frac{u^4}{\log u}\right)^{1/5},\qquad c>0,
\end{equation}
then the total-variation distance from independent unit-rate gamma
coordinates has the explicit limit
\begin{equation}\label{mc:eq:criticalTV}
 2\Phi\left(\frac12\sqrt{\frac{2\theta c^5}{25\pi}}\right)-1.
\end{equation}
If $h$ is smaller than the scale in \eqref{mc:eq:criticalheight} by a
factor tending to zero, that distance tends to zero.  If it is larger
by a factor tending to infinity, the distance tends to one, with no
additional restriction on $h/u$.  Indeed, in the latter case restrict
to the subblock with cutoff
\[
 h'=\min(h,u/\log u).
\]
Then $h'=o(u)$ and $h'/(u^4/\log u)^{1/5}\to\infty$, so
Theorem~\ref{mc:thm:threshold} gives distance tending to one for the
subblock; marginalization gives the same lower bound for the full
block.  Thus \eqref{mc:eq:criticalTV}, with its continuous endpoint
values zero and one at $c=0$ and $c=\infty$, describes every limiting
ratio $h/(u^4/\log u)^{1/5}\to c\in[0,\infty]$.

The fifth power and the logarithmic correction are
forced by the accumulated rate discrepancy across the growing block,
not by a bound on any single coordinate.  This threshold is smaller
than the scale $h\asymp u$ in Theorem~\ref{mc:thm:macroscopic}, where
the conditioning creates a nonvanishing dependence even after the
correct unequal rates are retained.

Under RH, the paired Hadamard product supplies the arithmetic identity
\begin{equation}\label{mc:eq:xilaw}
 \mathbb E e^{-tX}
 =\left(\frac{\xi(1/2)}{\xi(1/2+\sqrt t)}\right)^\theta.
\end{equation}
This is the precise bridge from the supplied article's gamma law to
the conditional theorems proved here.  The general probability theorem is
independent of RH; identification \eqref{mc:eq:xilaw} is where RH
enters.  It makes no assertion about simple zeros, zero correlations,
or unproved linear independence of ordinates.

\subsection{Complex characters and finite products}

The construction also covers a primitive nonreal Dirichlet character
$\chi$, by pairing it with its conjugate.  With its usual completed
function $\Lambda_\chi$ and central order
$m=\operatorname{ord}_{s=1/2}\Lambda_\chi(s)$, put
\begin{equation}\label{mc:eq:complexproduct}
 H_\chi(w)=
 \frac{\Lambda_\chi(1/2+w)\Lambda_{\bar\chi}(1/2+w)}{w^{2m}},
\end{equation}
with the removable singularity at zero filled in.  The functional
equations make $H_\chi$ even.  On the real axis the numerator is
$|\Lambda_\chi(1/2+w)|^2$, so
\[
 H_\chi(0)=\left|\frac{\Lambda_\chi^{(m)}(1/2)}{m!}\right|^2>0.
\]
Under Dirichlet GRH, the paired Hadamard factorization gives
\[
 \frac{H_\chi(w)}{H_\chi(0)}
 =\prod_{\gamma_j>0}\left(1+\frac{w^2}{\gamma_j^2}\right),
\]
where the positive zeros of the two completed functions are combined,
retaining their multiplicities.  The product converges locally
uniformly because $\sum_j\gamma_j^{-2}<\infty$.  Evenness eliminates
the linear exponential factor in the order-one Hadamard product.
Thus the gamma sum has Laplace transform
\[
 \mathbb E e^{-tX}
 =\left(\frac{H_\chi(0)}{H_\chi(\sqrt t)}\right)^\theta.
\]
Its positive spectral count satisfies
$N(T)\sim2T\log T/(2\pi)$.

More generally, multiply finitely many such conjugate-pair factors,
real primitive Dirichlet factors with their central powers removed,
and copies of $\xi(1/2+w)$.  If $d$ is the number of degree-one
completed factors before conjugate pairing, the positive spectrum has
\[
 N(T)\sim\frac{d}{2\pi}T\log T.
\]
The profiles $A$ and $V$ in \eqref{mc:eq:riemannprofiles} are unchanged.
The first two equivalents in \eqref{mc:eq:riemannscales} acquire a
factor $d$, and the threshold parameter becomes
\[
 \tau_u^2\sim\frac{\theta d\,h^5\log h}{10\pi u^4}
 \qquad(h\to\infty).
\]
At the critical height \eqref{mc:eq:criticalheight}, the limiting
distance from independent unit-rate gamma coordinates is therefore
\begin{equation}\label{mc:eq:degreeprofile}
 2\Phi\left(c^{5/2}\sqrt{\frac{\theta d}{50\pi}}\right)-1.
\end{equation}
The leading transition depends only on the total degree and the gamma
shape.  The conductors, root numbers, locations of individual zeros,
and their multiplicities affect the exact saddle and finite-scale
corrections, but not this limiting profile.

\begin{corollary}[The sharp exponent in the number of coordinates]
\label{mc:cor:four-fifths}
For any fixed degree $d$ in the preceding construction, write
$n=N(u)$ and $k=N(h)$.  If
\[
 \frac{k}{n^{4/5}}\longrightarrow\kappa\in[0,\infty],
\]
then the total-variation distance from $k$ independent unit-rate gamma
coordinates tends to
\begin{equation}\label{mc:eq:countprofile}
 2\Phi\left(\frac12\sqrt{\frac{125\theta}{256}\,\kappa^5}\right)-1,
\end{equation}
with value one at $\kappa=\infty$.  In particular, unit-rate
equilibration holds for $k=o(n^{4/5})$, and fails maximally when
$k/n^{4/5}\to\infty$.  Independence with the correct unequal rates
holds on the larger range $k=o(n)$.
\end{corollary}

\begin{proof}
At $h=c(u^4/\log u)^{1/5}$, the counting asymptotic gives
\[
 \frac{N(h)}{N(u)^{4/5}}
 \longrightarrow\frac{4c}{5}
                     \left(\frac d{2\pi}\right)^{1/5}.
\]
Substitute
$c=(5\kappa/4)(2\pi/d)^{1/5}$ into
\eqref{mc:eq:degreeprofile}; the degree cancels and
$\tau^2=125\theta\kappa^5/256$.
Monotonicity of the count brackets any sequence with a finite positive
limiting ratio between such critical cutoffs.  The zero and infinity
cases follow by the same bracketing and marginalization.  Finally,
regular variation implies $N(h)/N(u)\to0$ if and only if $h/u\to0$,
so Theorem~\ref{mc:thm:macroscopic} gives the last assertion.
\end{proof}

% Suggested bibliography keys/entries for the parent report:
% Article: pinned ProveIt Article/sections/ana.tex, blob
% 6b414e43f44bd89c5a7fce0c0f2a0cf08626cc8b, specifically
% ana:thm:gammasum and ana:thm:smalldev, lines 2955--3057.
% CampenhoutCover1981: Jan M. van Campenhout and Thomas M. Cover,
% Maximum Entropy and Conditional Probability, IEEE Trans. Inform. Theory
% IT-27 (1981), no.4, 483--489.
% https://isl.stanford.edu/people/cover/papers/transIT/0483camp.pdf
% MayerWolfZeitouni1993: Eddy Mayer-Wolf and Ofer Zeitouni,
% The Probability of Small Gaussian Ellipsoids and Associated Conditional
% Moments, Ann. Probab. 21 (1993), no.1, 14--24.
% https://doi.org/10.1214/aop/1176989391
% Author preprint: https://www.wisdom.weizmann.ac.il/~zeitouni/pdf/eddyelip.pdf
